\documentclass[pdftex,12pt,twoside,footsepline,headsepline,a4paper,final]{scrreprt}
\usepackage[verbose,a4paper,bottom=25mm,top=25mm]{geometry} %showframe,
\usepackage[T1]{fontenc}
\usepackage[utf8]{inputenc}
\usepackage{ngerman}
\usepackage{color}

\newcommand{\platz}{\vspace{2.1cm}}


\begin{document}

\begin{huge}
\textbf{Projektpraktikum 2026-2027}
\end{huge}

\par\noindent\rule{\textwidth}{0.4pt}

$_{}$



\begin{Large}
\textcolor[rgb]{0.31,0.51,0.74}{\emph{Erwartungen an das Projektpraktikum}}
\end{Large}


$_{}$


Die folgenden sechs Fragen sollen Ihnen helfen, Ihre Motivation und Ihre Vorkenntnisse zu bewerten.
Sie k\"{o}nnen Teile der Gesamtreflexion wie auch der Einleitung mit Hilfe Ihrer Antworten formulieren.

\begin{enumerate}

\item Ich nehme am Projektpraktikum teil, weil \ldots

\platz

\item Ich bringe folgende mathematische, physikalische und sonstige Vorkenntnisse und / oder
praktische Erfahrungen mit \ldots

\platz


\item Das Projektpraktikum wird f\"{u}r mich erfolgreich sein, wenn Folgendes eintritt \ldots

\platz


\item Folgendes darf hier nicht passieren \ldots

\platz


\item Das ist mir sonst noch wichtig \ldots

\platz


\item Ich bin bereit, Verantwortung sowohl f\"{u}r das Gelingen des Projektpraktikums als auch f\"{u}r
meinen eigenen Lernprozess zu \"{u}bernehmen und werde folgenden Beitrag leisten \ldots


\platz

\end{enumerate}


Quelle: E. Karagiannakis, Universit\"{a}t Ulm, Portfolio-Weiterbildung 2012

\end{document} 